\subsection{The Emergence of Modern Hopfield Networks}

The limitations of classical Hopfield networks motivated the development of more expressive associative memory architectures. A major advancement was introduced by Ramsauer et al. with the modern Hopfield network, which replaces discrete attractor states with memories represented in continuous high-dimensional spaces \cite{ramsauer2021hopfield}.

Instead of storing memories implicitly within fixed network parameters, modern Hopfield networks represent memories explicitly as continuous vectors. Retrieval is performed by comparing an input query representation against stored memory representations, allowing the system to identify relevant memories based on similarity.

An important result of this formulation is its connection to the attention mechanism used in Transformer models. Ramsauer et al. showed that the update rule of modern Hopfield networks is mathematically equivalent to the attention mechanism, establishing attention as a form of continuous associative memory retrieval \cite{ramsauer2021hopfield,vaswani2017attention}.

\subsection{External Memory Representation}

A key distinction between classical and modern Hopfield networks is the separation between model parameters and stored memories. Classical Hopfield networks encode memories within the connection weights of the network, meaning that adding new memories requires modifying the model itself. 

Modern Hopfield networks instead maintain memories as an external set of representations. This enables memories to be added or updated independently of the underlying retrieval mechanism, avoiding the need to retrain the model when new information becomes available.

This property makes modern Hopfield networks particularly relevant for adaptive machine learning systems, where the stored knowledge may change over time while the core model remains unchanged.

\section{Associative Memory}

Associative memory refers to the ability of a system to store and retrieve information based on relationships between patterns. In artificial intelligence (AI), associative memory models are designed to recover stored information from partial or incomplete inputs by learning associations between representations. Early approaches, such as Hopfield Networks, modelled memories as stable states within interconnected networks, where the structure of connections influenced memory storage and retrieval \cite{hopfield1982}. Modern Hopfield Networks (MHNs) build upon this foundation by introducing high-dimensional representations and attention-based retrieval mechanisms, allowing associative memory models to achieve greater scalability and improved retrieval capabilities \cite{ramsauer2021hopfield}. The concept of associative memory is also observed in biological neural systems, where it supports processes such as learning, recognition, and recall.

\section{Transformers and Attention Mechanisms}

The Transformer architecture introduced a shift from recurrent sequence modelling
towards attention-based representation learning. Unlike recurrent architectures,
Transformers allow each input representation to directly attend to other
representations within a sequence, enabling the model to dynamically capture
relationships between different elements~\cite{vaswani2017attention}.

The core operation of Transformers is the attention mechanism:

\[
Attention(Q,K,V)=
\operatorname{softmax}
\left(
\frac{QK^T}{\sqrt{d_k}}
\right)V
\]

where queries ($Q$) retrieve relevant information from keys ($K$) and combine
the corresponding values ($V$). This mechanism allows models to retrieve
relevant representations based on learned similarity relationships.

This behaviour is closely related to Modern Hopfield Networks. Ramsauer et al.
showed that the update rule of Modern Hopfield Networks can be expressed as an
attention operation~\cite{ramsauer2021hopfield}. The associative retrieval
process of an MHN can be represented as:

\[
\boldsymbol{\xi}^{new}
=
X\operatorname{softmax}
(\beta X^T\boldsymbol{\xi})
\]

where the query representation retrieves relevant stored patterns from a
continuous memory space. This formulation is equivalent to the attention
operation, demonstrating that attention mechanisms can be interpreted as a form
of scalable associative memory retrieval.

Due to the extensive ecosystem and practical success of transformer-based
architectures, this thesis focuses on transformer models rather than standalone
Modern Hopfield Networks. However, the retrieval behaviour of these models is
analysed from the perspective of associative memory. Transformer-based
bi-encoders and cross-encoders can therefore be viewed as practical
implementations of associative retrieval, where representations are either
retrieved through embedding similarity or compared through direct attention-based
interaction.

\subsection{Limitations of Attention-Based Associative Retrieval}

Although attention mechanisms provide an effective method for associative
retrieval, they introduce computational challenges when applied to large-scale
datasets. The standard self-attention operation requires computing pairwise
similarities between input representations, resulting in a computational
complexity of $O(n^2)$ with respect to the sequence length. This quadratic
scaling limits the applicability of attention mechanisms for very long inputs or
large retrieval collections.

Similarly, while Modern Hopfield Networks provide a scalable formulation of
associative memory compared to classical Hopfield Networks, retrieval still
requires computing similarities between queries and stored representations.
When the number of stored patterns becomes large, the cost of comparing against
the entire memory space becomes computationally expensive.

These limitations motivate the use of efficient retrieval architectures that
reduce the number of comparisons required during inference. In this thesis, the
two-stage retrieval architecture addresses this challenge by using a
bi-encoder to efficiently retrieve a limited set of candidate entities from a
large regulatory corpus, followed by a cross-encoder that performs more
computationally expensive attention-based interaction modelling on the reduced
candidate set.